%% kapa: exam=1 round=34 year=2023 period=2 form=A subject=accounting subject_ko=회계학 seq=34 number=74
%% source: Q-Net 감정평가사 기출문제 (cst00309) artlSeq=5211789 file=raw/1008/5211789/2023년도 제34회 감정평가사 1차시험 2교시.pdf
%% conversion: kordoc 4.24.51 -> tools/gen1.py (auto); review=proofread-sonnet; flags=none
\begin{question}{74}
(주)감평은 제품 생산에 필요한 부품 400단위를 매년 외부에서 단위당 ￦1,000에 구입하였다. 그러나 최근 외부구입가격 인상이 예상됨에 따라 해당 부품을 자가제조하는 방안을 검토하고 있다. 다음은 (주)감평이 부품 100단위를 자가제조할 경우의 예상제조원가 자료이다.
\par{\small\centering\begin{tabular}{|l|l|}\hline
직접재료원가 & ￦25,000 \\ \hline
직접노무원가 & \hspace{1em}30,000 (￦100/직접노무시간) \\ \hline
변동제조간접원가 & \hspace{1em}20,000 (직접노무원가의 2/3) \\ \hline
고정제조간접원가 & \hspace{1em}100,000 (전액 유휴생산설비 감가상각비) \\ \hline
\end{tabular}\par}
(주)감평은 현재 보유하고 있는 유휴생산설비를 이용하여 매년 필요로 하는 부품 400단위를 충분히 자가제조할 수 있을 것으로 예상하고 있으며, 부품은 한 묶음의 크기를 100단위로 하는 묶음생산방식으로 생산할 예정이다. 해당 부품을 자가제조하는 경우, 직접노무시간이 학습률 90\%의 누적평균시간 학습모형을 따를 것으로 추정된다. (주)감평이 부품 400단위를 자가제조할 경우, 단위당 제조원가는?
\begin{choices}
  \item ￦655
  \item ￦712
  \item ￦750
  \item ￦905
  \item ￦1,000
\end{choices}
\end{question}
\answer{4}
